\section{LLM-Assisted Assembly Transformation}\label{sec:llm}
\subsection{Task and Supported Transformation Families}
The transformation path asks an LLM to produce a complete modified
assembly representation for a selected family: EQ, NOP insertion, BCF,
or IBP. These families retain their meanings from the native engine;
the use of a model changes how the proposed representation is produced.
Each request selects one family so that the generated artifact and its
validation policy have an identifiable experimental condition.

PII and assembly wrapping remain native paths in the version examined
here. They are not described as LLM-generated techniques. Likewise, a
native condition combining transformations is not automatically
equivalent to several model calls in sequence. The experiment identifies
the implemented path rather than grouping all uses of the same family
under one label.

\subsection{Assembly Representation and Data Boundary}
The engine prepares an assembly representation with symbolic labels
for the relevant locations and preserves the data representation
outside the model-generated instruction text. This preparation defines
an intermediate input to the transformation task. It also determines
the reference point for several structural checks and similarity
measurements.

The model's text is assembled and associated with the data retained by
the engine. The output is therefore a product of model generation,
assembly, data handling, and validation. The experimental record retains
the identity of the original file, the intermediate representation,
and the final bytes, so that preparation-induced changes can be
distinguished from changes in the model's output.

\subsection{Zero-Shot and Two-Shot Conditions}
The initial zero-shot request contains the task specification without
worked examples. The initial two-shot request contains two examples
in addition to that specification. The example identities and prompt
versions are part of the experimental condition. Differences in their
content, not only their number, can influence the comparison.

These labels apply to the first request in a run. After rejection, a
later request may contain the previous answer and validation feedback.
Such a run remains assigned to its initial prompting condition, but its
later attempts are explicitly feedback-assisted. Reporting only the
final accepted output would hide this additional interaction.

The experimental unit is a run under a specified input, transformation,
prompt condition, and generation configuration. Repeated run numbers
are not claimed to be independently controlled model seeds unless the
service accepts and records such a parameter. Temperature, token
budget, model identifier, and service configuration are recorded
separately from the run index.

\subsection{Validation and Feedback}
Validation covers several kinds of structural evidence. Assembly must
succeed, the produced code must meet the decoding checks, tracked branch
destinations must have appropriate boundaries, and the expected data
representation must remain consistent. Additional checks depend on the
selected family and examine whether the output still contains the
required instruction structure or applies a nontrivial transformation.

These checks are deliberately narrower than semantic equivalence.
An output can pass structural validation yet behave differently; a
conservative check can also reject an output that would have performed
the expected behavior. The assessment records the reason for each
rejection rather than treating all rejected answers as one category.

For supported test cases, a functional verifier can participate in the
generation loop. If that verifier is also used to select accepted
outputs, a later test based on the same observable is conditional on
selection. The article distinguishes this pipeline acceptance criterion
from an independent functional assessment and reports attempts that
exhaust the allowed budget.

\subsection{What Is Attributed to the Model}
Three aspects are considered separately: whether an answer can be
assembled, whether the complete pipeline accepts it, and whether an
accepted output satisfies the reported functional check. Deterministic
preparation and validation are part of the system contribution rather
than capabilities attributed to the model alone.

The run record includes elapsed time, API usage when supplied by the
service, rejected attempts, and the accepted artifact. Estimated cost
is derived from the stated tariff snapshot; it is not presented as an
invoice. Resumed jobs and reused outputs are identified so that they
are not counted as additional generation calls.
